\newcommand{\DoNotLoadEpstopdf}{}
\documentclass[11pt,reqno]{amsart}
\usepackage[T1]{fontenc}
\usepackage{lmodern,amsmath,amssymb,amsthm,mathtools,mathrsfs,microtype}
\usepackage[margin=1.1in]{geometry}
\usepackage[hidelinks]{hyperref}
\hypersetup{pdftitle={Collision records, arithmetic, and raw local inversion in a triangular Lorentz gas},pdfauthor={Qian Qi}}
\setlength{\emergencystretch}{2em}
\allowdisplaybreaks[2]
\numberwithin{equation}{section}
\newtheorem{theorem}{Theorem}[section]
\newtheorem{lemma}[theorem]{Lemma}
\newtheorem{proposition}[theorem]{Proposition}
\newtheorem{corollary}[theorem]{Corollary}
\theoremstyle{definition}
\newtheorem{definition}[theorem]{Definition}
\newtheorem{inputstatement}[theorem]{Input statement}
\theoremstyle{remark}\newtheorem{remark}[theorem]{Remark}
\newcommand{\R}{\mathbb R}\newcommand{\Z}{\mathbb Z}
\newcommand{\one}{\mathbf1}\newcommand{\dd}{\,\mathrm d}
\newcommand{\esssup}{\operatorname*{ess\,sup}}
\newcommand{\diver}{\operatorname{div}}\newcommand{\TV}{\operatorname{TV}}
\title[Collision records and raw local inversion]{Collision records, arithmetic, and raw local inversion in a triangular Lorentz gas}
\author{Qian Qi}
\date{October 12, 2026. A2-DYN, revision 93}
\subjclass[2020]{37D50, 37A50, 60F05, 28A75}
\keywords{Lorentz gas, exact collision records, optical traces, signed source balance, one-sided persistence, pointwise local limits}
\begin{document}
\raggedbottom
\begin{abstract}
We return to the upper pointwise local limit problem for the exact
collision record of actual returns in a triangular Lorentz gas.
The positive optical denominator cancels each individual incidence
singularity in the absolute angular density, giving a collision-uniform
one-half modulus on every physical sheet. Its complete fixed-count
signed balance shows that the original critical-or-short-residence
selector creates only negative interior cuts. Positive increments come
from actual entrance traces and positive physical density growth.
We prove a prescribed-scale pointwise comparison retaining only the
ancestors of the selected target fiber, together with an upper-spike
persistence dichotomy on that same source. The complete noncritical
growth term has ordered height of order the roof width, with a leading
constant independent of every fixed positive gradient cutoff. Actual
ancestral entrances and critical positive growth remain explicit in
the scalar upper bound. No vanishing of their concentration is inferred,
and the original unrestricted pointwise theorem remains to be completed.
All earlier finite-norm and conditioning results are retained.
\end{abstract}
\maketitle
\input{journal-problem-v93}

\subsection{The upper sign and its physical origins}
The organizing assertion is still \eqref{eq:v76-target}. The historical
one-sided lower law and the finite-norm conclusions remain available,
but do not answer the upper-spike question. The new argument studies
the original positive weighted density before taking its absolute
roof increments. It uses the monotonicity of the actual exceptional
selector along each maximal physical segment, not an embedding of a
finite norm into an essential supremum.

Theorems~\ref{thm:v93-balance} and \ref{thm:v93-pointwise} give its
signed balance and a pointwise inequality at a prescribed roof radius.
Theorem~\ref{thm:v93-gain} estimates the full noncritical positive
variation after the complete physical summation. Thus a high value
either persists on the prescribed backward interval or is paid by a
quantified entrance or critical-growth contribution of the same target's
ancestors. The final upper bound in Theorem~\ref{thm:v93-upper-endpoint}
keeps those contributions rather than declaring them small.

\input{journal-inputs-v93}
\input{core/220_absolute_optical_traces}
\input{core/221_signed_prefix_birth_balance}
\input{core/222_upper_birth_and_critical_gain}

\section{The original endpoint and the retained consequences}
The comparison just proved acts at almost every original roof, not
only after integrating starting roofs. It also retains all true
entrances and losses. An interior cutoff contributes a negative trace,
so the original exceptional union does not acquire an uncontrolled
positive cutoff-boundary term. These conclusions do not use the
positive-reference denominator or suppress upper amplitudes by taking
a reciprocal.

What is not proved is that the two positive forcing terms in
\eqref{eq:v93-endpoint} vanish. A uniform one-sheet modulus cannot be
summed with an unestimated angular multiplicity. A critical point's
zero area does not remove its possible directional entrance trace.
The remaining task is therefore a complete-source estimate of these
specific positive mechanisms or another source-specific upper bound.
No new unrestricted charge-zero or pointwise likelihood assertion is
made here.

All earlier weak, strong subcritical, logarithmic, diffusive-window,
one-sided lower, reciprocal and conditioning theorems retain exactly
their previous hypotheses. Both revision-88 variants remain distinct.
The technical archive includes the full preceding article and every
inherited proof; no alternate model, averaged target or pooled label
has replaced \eqref{eq:v76-target}.

\appendix
\input{appendices/v93_dependency_map}
\input{journal-references-v93}
\end{document}
